%!TEX root =  tfg.tex
\chapter{Contexto}

\begin{quotation}[Novelist]{Ernest Hemingway (1899--1961)}
The good parts of a book may be only something a writer is lucky enough to overhear or it may be the wreck of his whole damn life -- and one is as good as the other.
\end{quotation}

\begin{abstract}
Este capítulo detalla los motivos que impulsan el desarrollo de Via Inferni, analizando la oportunidad de mercado en el género roguelike y el público al que va dirigido. Asimismo, se definen de forma precisa los objetivos técnicos y funcionales que guiarán la construcción del videojuego bajo una metodología de Ingeniería del Software.\end{abstract}

\section{Motivación}
El desarrollo de este proyecto nace de la intersección entre la literatura clásica y las tendencias actuales de la industria del videojuego independiente. El género roguelike ha demostrado una gran viabilidad comercial debido a su alta rejugabilidad y costes de producción contenidos. Sin embargo, existe una oportunidad para integrar narrativas profundas y estructuras literarias complejas, como la Divina Comedia, en mecánicas de juego procedurales.

La hipótesis es que la estructura de los nueve círculos del infierno es ideal para un diseño de niveles ascendentes en dificultad. El público objetivo son jugadores aficionados a retos de habilidad y la estética pixel art, que buscan experiencias de juego cortas pero intensas.

\section{Listado de objetivos}

\begin{description}
\item \textbf{Objetivo 1. Diseño y desarrollo de una arquitectura de juego modular.} Implementar un sistema sólido en Unity que permita la gestión de entidades (jugador, enemigos, ítems) y la separación de la lógica de juego de la representación visual.
\item \textbf{Objetivo 2. Generación procedural de niveles.} Crear un algoritmo de generación aleatoria de salas que garantice que cada partida sea única, respetando la coherencia temática de los nueve círculos del infierno.
\item \textbf{Objetivo 3. Implementación del sistema de "Cambio Dual".} Desarrollar la mecánica que permite alternar entre dos personajes (Dante y Virgilio), cada uno con estadísticas y habilidades diferenciadas, aportando una capa estratégica al combate.
\item \textbf{Objetivo 4. Creación de un ecosistema de enemigos y jefes.} Programar una inteligencia artificial variada con al menos 15 tipos de enemigos y jefes finales con patrones de ataque específicos (como Cancerbero y Satán).
\item \textbf{Objetivo 5. Gestión de progresión y economía.} Diseñar e implementar un sistema de inventario, objetos pasivos/consumibles y una tienda que permita al jugador mejorar sus capacidades durante la partida.
\end{description}